\noindent\textbf{Blatt 0}

\erg{1}{a) $\frac{2}{5}$ \quad b) $\frac{2}{3}$ \quad c) $\frac{3}{8}$ \quad d) $\frac{35}{100}$, Zähler 35 \quad e) $\frac{7}{8}$}
\erg{2}{a) $0{,}5$ \quad b) $0{,}25$ \quad c) $\frac{2}{5}$ \quad d) $0{,}6$ \quad e) $0{,}125$ \quad f) $\frac{7}{100}$, Zähler 7}
\erg{3}{a) $3{,}2$ \quad b) $7{,}7$ \quad c) $12{,}3$ \quad d) $10{,}0$ \quad e) $5{,}6$ \quad f) $28{,}0$}
\erg{4}{a) $4$ \quad b) $0{,}35$ \quad c) $0{,}075$ \quad d) $24$ \quad e) $2{,}4$ \quad f) $60$}
\erg{5}{a) $:5$, 1 Anhänger $0{,}80\,€$; $\cdot 8$: $6{,}40\,€$ \quad b) $:6$, 1 Brötchen $0{,}35\,€$; $\cdot 4$: $1{,}40\,€$ \quad c) $:3$, 1\,m $2{,}50\,€$; $\cdot 7$: $17{,}50\,€$ \quad d) $:4$, 1 Los $0{,}75\,€$; $\cdot 16$: 16 Lose}
\erg{6}{a) $40\,€$ \quad b) $15\,$kg \quad c) $18\,$m \quad d) $18\,€$ \quad e) $0{,}70\,€$}
\erg{7}{a) Mia hat Zähler und Nenner einfach hintereinander hinter das Komma geschrieben, statt zu teilen; richtig: $3 : 8 = 0{,}375$ \quad b) $0{,}45$}

\par\medskip\noindent\textbf{Wofür Blatt 0 gebraucht wird:}
1 → Einheit 1 und 2 \quad 2 → Einheit 1 und 3 \quad 3 → Einheit 2 und 5 \quad 4 → Einheit 2 bis 4 \quad 5 → Einheit 2 bis 4 \quad 6 → Einheit 3 \quad 7 → Einheit 1 und 2
